\section{Introduction}
\label{cha:introduction}

The digital revolution has fundamentally transformed the way we create, consume, and interact with visual information.
Images are no longer only physical items around us, they are inherent components of our online experience, shaping perceptions and influencing decisions.
As the Internet continues to evolve, so too grows the demand for rich and engaging visual content.
This demand, however, comes at a cost. Images, while visually appealing, are significantly impacting website performance, not only by generating the highest number of requests, but also by consuming the most bytes during page load of an average site.
Existing studies have shown a strong correlation between page load times and user engagement, conversion rates, and even search engine rankings.
Minimizing image size is therefore fundamental for efficient bandwidth use and plays a crucial role in optimizing website performance.

Image optimization practices for the web have evolved over time, following developments of the state-of-the-art image formats and codecs.
These solutions perform operations such as image resizing and format conversion, and may allow for storing optimized versions of original files for future use.
While the optimization itself is a relatively specialized task, its deployment models follow familiar and established patterns in the general software industry: we can find them in software-as-a-service offerings, integrated into web hosting and content delivery platforms, but also as open-source solutions for self-hosting.

This last option, leveraging open-source image processing libraries, presents a compelling opportunity for our research.
Unlike proprietary solutions, open-source projects offer the flexibility to be deployed across a variety of environments, granting builders granular control over the optimization process.
This allows not only to measure influence of different libraries and optimization parameters, but also – and crucially for this thesis – enables realistic benchmarking of the underlying platforms on which these operations are performed.
By employing these workloads, our research aims to accurately measure cloud service performance under the realistic conditions of real-world image optimization scenarios.

Given the event-driven nature of the optimization workload in response to requests for images, serverless cloud compute emerges as a well-suited deployment option for these systems.
However, the varying demands for image processing, including different image sizes and diverse range of image formats present on the web, raise questions about the suitability and performance consistency of these platforms for the use-case.
Furthermore, the emergence of function as a service (FaaS) platforms introduces a new dimension to cloud-based, self-hosted image optimization.
FaaS offers a compelling execution environment for image processing tasks thanks to its on-demand scalability and event-driven nature.
However, the performance variability inherent in FaaS platforms, attributed to diverse factors, may impact efficiency of the process and applicability of the environment for the use case.


\subsection{Objectives of the Thesis}
This thesis aims to investigate whether FaaS performance variability has measurable and meaningful influence for image optimization workloads.
We will define set of representative metrics related to image optimization, reflecting realistic web usage.
We will design and implement the optimization workflows for selected serverless environments.
We will design and implement scheduling mechanism allowing us to simulate requests for optimized images.
We will collect performance metrics across all parameters, apply statistical methods for result comparisons, and visualize the results.
We will investigate if more intensive optimization can be cost-efficient, considering optimization result caching, and content reuse for multiple viewers.
We will provide recommendations that are within the scope of the benchmark, and answer the primary research question in the context of the scenario. 


\subsection{Structure}
Provide a roadmap for the reader, outlining the organization of the thesis and briefly summarizing the content of each chapter.
